\section{Актуальность разработки модуля рекомендаций параметров тренировок}

Вестибулярная реабилитация требует длительных и регулярных занятий, а положительный эффект
формируется за счёт прогрессии нагрузки. На приёме
лечащий врач определяет тактику ведения пациента~--- какие упражнения допустимы, с какой
ориентировочной сложностью начинать, как часто необходим контроль состояния. Между визитами
основная часть тренировок выполняется в домашних условиях, без ежедневного участия
специалиста.

Разрабатываемый тренажёр ориентирован именно на такой режим: пациент выполняет упражнения
перед камерой, программный комплекс фиксирует объективные результаты и сохраняет их в базе данных. Наряду с
контролем техники выполнения и ведением истории тренировок возникает необходимость
использовать накопленные данные для адаптации нагрузки пациента между консультациями с врачом. Исходя из этого, было принято решение реализовать модуль рекомендаций параметров тренировок.

Актуальность его разработки обусловлена двумя факторами.

\textbf{Первый фактор} -- разрыв по времени между назначением врача и ежедневной практикой. Пациенту целесообразно заниматься несколько раз в неделю, тогда как очный контакт с врачом возможен существенно реже. За этот период меняется переносимость нагрузки: после перерыва результаты могут снижаться, при устойчивом выполнении -- улучшаться. Запись в плане лечения с фиксированными числовыми параметрами на весь межвизитный интервал не отражает такую динамику.

\textbf{Второй фактор} -- риск неоптимальной самостоятельной настройки. Пациент может самостоятельно изменить параметры тренировки: как усложнить, так и облегчить себе задание. Однако при этом он, как правило, не учитывает полную историю своих предыдущих занятий -- например, не видит устойчивых трендов успешности или частоты ошибок. В результате возможны два нежелательных сценария. При чрезмерном усложнении нагрузка становится неадекватно высокой, что ведёт к ухудшению техники выполнения упражнений, росту числа ошибок, снижению мотивации и, в конечном счёте, к пропуску тренировок. При неоправданном облегчении эффект реабилитации снижается, а пациент теряет стимул для прогресса.

Важно определить границы применения модуля. Рекомендательная подсистема \emph{не заменяет}
лечащего врача и \emph{не назначает} новые упражнения или этапы реабилитации. Она опирается
на \textbf{базовые параметры, зафиксированные по назначению врача} для данного пациента и
типа упражнения, и предлагает их уточнение на текущую сессию с учётом последних
результатов~--- в пределах ограниченного шага изменения сложности. Врач задаёт план и
опорные значения; между визитами система помогает применять этот план с учётом
фактической динамики тренировок.

Практическая значимость состоит в том, что пациент перед каждым занятием получает понятную рекомендацию по параметрам нагрузки. Это позволяет ему тренироваться эффективно без постоянного контроля врача. 

\subsection{Постановка задачи модуля рекомендаций параметров тренировок}

В программном комплексе девять упражнений. Перед запуском каждого
из них задаётся свой набор параметров:

\begin{itemize}
    \item \textbf{Упражнение 1:} количество объектов, время на объект,
    фон, звук;
    \item \textbf{Упражнение 2:} количество объектов, время на объект,
    фон;
    \item \textbf{Упражнение 3:} количество объектов, скорость, фон;
    \item \textbf{Упражнение 4:} количество объектов, время на объект,
    скорость, фон;
    \item \textbf{Упражнение 5:} количество объектов, время на объект,
    скорость, фон;
    \item \textbf{Упражнение 6:} количество объектов, время на объект,
    фон;
    \item \textbf{Упражнение 7:} количество объектов, время на объект,
    диапазон движений шеи, фон;
    \item \textbf{Упражнение 8:} количество объектов, интервал смены
    цвета, скорость, фон;
    \item \textbf{Упражнение 9:} количество объектов, интервал смены
    цвета, скорость, фон.
\end{itemize}

Для рекомендательной системы из этого полного набора используются
только \textbf{нагрузочные} параметры~--- те, что непосредственно
задают сложность упражнения:

\begin{itemize}
    \item количество объектов;
    \item время на объект;
    \item скорость;
    \item интервал смены цвета;
    \item диапазон движений шеи.
\end{itemize}

Параметры \textbf{фона} и \textbf{звука} в рекомендательную систему
не входят. Их назначает только врач: они зависят от диагноза,
переносимости отвлекающих стимулов и индивидуальных особенностей
пациента и не выводятся из результатов тренировок.

Отдельно хранится \textbf{базовое назначение врача}~--- набор параметров
(нагрузочных и сенсорных), который врач задаёт пациенту для
самостоятельных занятий между визитами. Именно назначение врача
является опорной точкой, от которой система предлагает корректировку
нагрузочных параметров на очередное занятие.

После каждого упражнения в базу данных сохраняется запись, содержащая:

\begin{itemize}
    \item идентификатор пользователя;
    \item дату и время тренировки;
    \item количество объектов в сессии;
    \item количество пойманных объектов;
    \item фон;
    \item коэффициент сложности;
    \item итоговый балл;
    \item нагрузочные параметры, с которыми проводилась сессия
    (время на объект, скорость, интервал смены цвета или диапазон шеи~---
    в зависимости от упражнения).
\end{itemize}

Для работы рекомендательной системы из сохранённых данных
используются:

\begin{itemize}
    \item дата тренировки;
    \item количество объектов и количество пойманных~--- для расчёта
    успешности;
    \item коэффициент сложности и итоговый балл;
    \item нагрузочные параметры, с которыми выполнялась сессия.
\end{itemize}

Итого задача модуля рекомендаций состоит в следующем.

Для \textbf{нагрузочных параметров тренировок}
\begin{itemize}
    \item количество объектов,
    \item время на объект,
    \item скорость,
    \item интервал смены цвета,
    \item диапазон движений шеи,
\end{itemize}
используя
\begin{itemize}
    \item \textbf{базовое назначение врача} для данного пациента
    и упражнения,
    \item \textbf{данные из истории тренировок:} дату, количество
    объектов, количество пойманных, коэффициент сложности, итоговый
    балл и нагрузочные параметры последних сессий,
\end{itemize}
получить \textbf{рекомендуемые нагрузочные параметры} на очередное
занятие как корректировку назначения врача. Параметры фона и звука передаются
без изменений из назначения врача.

\subsection{Анализ существующих решений задачи адаптации параметров тренировки}

Задача модуля рекомендаций относится к классу \emph{адаптивного дозирования
нагрузки}: по результатам выполнения упражнений нужно изменить параметры
следующей сессии. В общем виде это отображение
\[
(\mathbf{B},\, \mathcal{H}) \;\longrightarrow\; \mathbf{P}^{*},
\]
где $\mathbf{B}$~--- назначение специалиста (опорная точка),
$\mathcal{H}$~--- история выполнения, $\mathbf{P}^{*}$~--- параметры
на очередную тренировку. Ниже рассмотрены подходы, описанные в литературе
для задач той же структуры, с указанием математической схемы каждого
решения.

\subsubsection{Системы мониторинга без автоматической адаптации}

В вестибулярной реабилитации уже существуют цифровые платформы, которые
фиксируют выполнение упражнений дома и передают данные специалисту.
В системе VestAid физиотерапевт через веб-портал задаёт параметры
домашней программы (дозировка, скорость движения головы, фон и др.),
а приложение на планшете контролирует соблюдение этих параметров и
точность выполнения~\cite{hovareshti2021}. Математически это соответствует
фиксированному назначению:
\[
\mathbf{P}^{*} = \mathbf{B},
\]
без использования $\mathcal{H}$ для автоматической корректировки.
Авторы отмечают, что доступ к объективным показателям помогает врачу
управлять прогрессией между визитами; автоматическая подстройка сложности
по результатам названа перспективным направлением.

Аналогичная логика реализована в цифровой платформе вестибулярной
реабилитации Meldrum и соавт.~\cite{meldrum2022}: на приёме физиотерапевт
назначает индивидуальную программу, при повторных визитах корректирует
её вручную. Объективный контроль движений дополняет субъективную оценку,
но решение о смене нагрузки остаётся за специалистом.

Такие системы решают задачу контроля выполнения, но не задачу
автоматической рекомендации параметров между визитами.

\subsubsection{Пороговые правила (rule-based)}

Наиболее близкая к поставленной задаче математическая схема~---
пороговое правило на метрике качества выполнения. В работе Noble и
соавт.~\cite{noble2021} при реабилитации после инсульта сложность
целевого задания $r_k$ обновляется только если ошибка выполнения на
предыдущем шаге $e_{k-1}$ не превышает порог $\varepsilon$:
\[
r_k =
\begin{cases}
r_{k-1} + \alpha r^{*}, & \text{если } \|e_{k-1}\|_2 \leq \varepsilon,\\
r_{k-1}, & \text{иначе.}
\end{cases}
\]
Здесь $r_k$~--- аналог нагрузочного параметра, $e_{k-1}$~--- отклонение
фактического результата от ожидаемого. Структура совпадает с поставленной
задачей: дискретное действие «усложнить / не менять» принимается по
измеренной метрике. Авторы сравнивают это правило с итеративным обучением
с обратной связью (ILC) и показывают, что пороговый подход проще, но
может давать слишком крупные шаги при пограничных значениях ошибки.

В роботизированной реабилитации при обнаружении утомления сопротивление
уменьшается по аналогичной схеме: если индикатор превышает порог,
нагрузка снижается на фиксированную величину~\cite{emg2020}. Клиническое
правило «Rule of 10» формализует ту же идею: прекратить или ослабить
нагрузку, когда интегральный показатель перенапряжения превышает
пороговое значение~\cite{rule10}.

Общая математическая форма таких решений:
\[
a_k = R(x_k;\, \theta),
\]
где $x_k$~--- измеренный показатель (ошибка, успешность, утомление),
$\theta$~--- пороги, $R$~--- кусочно-постоянная функция, $a_k$~---
дискретное действие над параметрами.

\subsubsection{Нечёткий вывод (fuzzy inference)}

G\'{o}mez-Portes и соавт.~\cite{gomez2021} формализуют адаптацию плана
реабилитации после инсульта как отображение показателей выполнения
в изменения параметров упражнений, изначально заданных терапевтом:
\[
\Delta \mathbf{p} = \mathrm{Defuzz}\bigl( F(\boldsymbol{\Phi}) \bigr),
\]
где $F$~--- набор нечётких правил вида «если успешность низкая, то
уменьшить число повторений / увеличить время». Терапевт задаёт
исходный план, система предлагает его модификацию по мере прогресса
пациента. Аналогичная схема применяется в системе поддержки решений
для домашней моторной реабилитации~\cite{castro2023}: на входе~---
процент выполнения и уровень успешности, на выходе~--- приращения
параметров упражнений. Нечёткая модель математически обобщает пороговые
правила, допуская промежуточные значения («умеренно низкая успешность»).

\subsubsection{Гибридные системы (правила и случаи)}

Li и соавт.~\cite{zhou2024} предложили систему поддержки решений для
роботизированной реабилитации верхних конечностей на основе гибрида
rule-based reasoning (RBR) и case-based reasoning (CBR). Правила
формируют план тренировок по клиническим протоколам, а подбор
конкретных параметров в ходе занятия уточняется по историческим
случаям других пациентов с использованием алгоритма Random Forest.
Математически это двухуровневая схема:
\[
\mathbf{P}^{*} = \mathrm{CBR}\bigl(\mathrm{RBR}(\mathbf{B}),\, \mathcal{H}\bigr),
\]
требующая базы случаев из десятков пациентов и этапа обучения модели.

\subsubsection{Модельная адаптация и обучение с подкреплением}

В подходе на основе модели моторного улучшения~\cite{mi2020} сложность
задания меняется по непрерывной оценке $\widehat{MI}_k$, вычисляемой
из кинематических показателей каждой сессии:
\[
\mathbf{P}^{*}_k = h(\widehat{MI}_k),
\]
где $h$~--- функция, откалиброванная на выборке пациентов. ILC-подход
из~\cite{noble2021} непрерывно модулирует шаг изменения сложности в
зависимости от ошибки, что математически богаче порогового правила,
но требует модели динамики выполнения.

Задача может быть записана и как поиск политики
$\pi: \mathcal{S} \to \mathcal{A}$, максимизирующей ожидаемый прогресс.
Для этого нужна большая выборка переходов «состояние~--- действие~---
результат»; при малом числе сессий на одного пациента обучение
политики неустойчиво.

\subsubsection{Сопоставление подходов и вывод}

Для задачи рекомендации параметров в разрабатываемом тренажёре
существенны четыре ограничения:
\begin{enumerate}
    \item мало наблюдений на одного пациента (единицы--десятки сессий);
    \item дискретное пространство параметров (целые числа, три уровня
    скорости);
    \item обязательная привязка к назначению врача $\mathbf{B}_e$,
    а не к параметрам последней сессии;
    \item сенсорные параметры (фон, звук) исключены из автоматической
    корректировки.
\end{enumerate}

Каждый из рассмотренных подходов оценён по этим ограничениям.

\textit{Только назначение врача}
(системы мониторинга~\cite{hovareshti2021,meldrum2022}).
История $\mathcal{H}$ не используется для автоматической корректировки:
$\mathbf{P}^{*} = \mathbf{B}$. Обучение на большой выборке не требуется,
привязка к назначению врача соблюдается, дискретные параметры
поддерживаются. Недостаток~--- задача адаптации между визитами
не решается, хотя объективные данные уже накапливаются.

\textit{Пороговые правила (rule-based)}
~\cite{noble2021,rule10,emg2020}.
История используется: решение $a_k = R(x_k;\,\theta)$ принимается по
метрикам выполнения. Большая обучающая выборка не нужна~--- пороги
$\theta$ задаются экспертно. Привязка к $\mathbf{B}_e$ реализуется
через оператор корректировки $\mathrm{Adj}(\mathbf{B},\,\delta)$.
Дискретные параметры и ограниченный шаг сложности ($\delta \in \{-1,0,+1\}$)
естественны для этого подхода.

\textit{Нечёткий вывод (fuzzy inference)}
~\cite{gomez2021,castro2023}.
По свойствам близок к пороговым правилам: история используется,
обучение на большой выборке не требуется, исходный план задаёт
терапевт ($\mathbf{B}_e$), параметры дискретны. Отличие~--- пороги
заменены функциями принадлежности; настройка сложнее, но структура
решения та же.

\textit{Гибрид RBR + CBR}
~\cite{zhou2024}.
История используется; для CBR-части нужна база случаев из многих
пациентов и обучение модели (Random Forest). Привязка к назначению
врача частичная: правила задают план, параметры уточняются по
похожим случаям. Для прототипа с малым числом пользователей
требования к данным избыточны.

\textit{Model-based и ILC}
~\cite{noble2021,mi2020}.
История используется; требуется калибровка модели улучшения или
динамики ошибки на выборке пациентов. Назначение врача как опорная
точка в явном виде не заложено~--- параметры подстраиваются
непрерывно от текущего состояния. Непрерывная природа ILC
плохо согласуется с дискретными настройками тренажёра (три
уровня скорости, целые секунды).

\textit{Машинное обучение и обучение с подкреплением.}
История используется; необходима большая выборка переходов
«состояние~--- действие~--- результат». Привязка к $\mathbf{B}_e$
не гарантирована без дополнительных ограничений в модели.
Интерпретируемость решения ниже, чем у rule-based и fuzzy-подходов.

Из проведённого анализа следует, что наиболее применимы к поставленной
задаче \textbf{пороговые правила} и их \textbf{нечёткое обобщение}:
оба подхода не требуют обучения на большой выборке, допускают привязку
к назначению врача и работают с дискретными параметрами. Системы
мониторинга в вестибулярной реабилитации~\cite{hovareshti2021,meldrum2022}
подтверждают актуальность сбора объективных данных, но не закрывают
задачу автоматической рекомендации между визитами. Гибридные,
model-based и ML-подходы избыточны для прототипа с малым объёмом
данных на одного пациента.

На основании этого выбора в следующем подразделе описывается
метод решения задачи рекомендации параметров.

\subsection{Метод решения задачи рекомендации параметров}

На основании анализа существующих подходов (подраздел выше) для
разрабатываемого тренажёра выбран \textbf{пороговый rule-based метод}.
Его математическая схема совпадает с описанной в~\cite{noble2021,gomez2021}
для задач адаптивного дозирования: по метрикам истории определяется
дискретное действие, которое применяется к назначению врача. Выбор
обоснован тем, что при малом числе сессий на пациента метод не требует
фазы обучения, согласуется с дискретной природой параметров тренажёра
и сохраняет назначение врача как опорную точку~--- как в нечёткой
системе G\'{o}mez-Portes~\cite{gomez2021}, где исходный план задаёт
терапевт, а система предлагает его модификацию.

\subsubsection{Математическая постановка}

Пусть зафиксированы пациент $u$ и упражнение $e$ ($e = 1, \ldots, 9$).

\textbf{Базовое назначение врача} для пары $(u, e)$:
\[
\mathbf{B}_e = (\mathbf{B}_e^{\mathrm{load}},\, \mathbf{B}_e^{\mathrm{sens}}),
\]
где $\mathbf{B}_e^{\mathrm{load}}$~--- нагрузочные параметры (количество
объектов, время, скорость, интервал смены цвета, диапазон движений шеи),
$\mathbf{B}_e^{\mathrm{sens}}$~--- сенсорные (фон, звук).

\textbf{История} последних $N$ сессий:
\[
\mathcal{H}_{u,e} = \bigl\{ (t_i,\, n_i,\, c_i,\, \kappa_i,\, \sigma_i) \bigr\}_{i=1}^{N},
\]
где $t_i$~--- дата, $n_i$~--- число объектов, $c_i$~--- число пойманных,
$\kappa_i$~--- коэффициент сложности, $\sigma_i$~--- итоговый балл.

Из истории вычисляется вектор признаков:
\[
\boldsymbol{\Phi}(\mathcal{H}_{u,e}) = (\bar{S},\, T,\, A),
\]
где
\[
\bar{S} = \frac{1}{N}\sum_{i=1}^{N}\frac{c_i}{n_i} \cdot 100\%
\]
--- средняя успешность;
$T$ --- разность средних успешностей последних и предыдущих сессий
(при $N \geq 6$: последние 3 минус предыдущие 3; при $N = 4, 5$:
последние 2 минус предыдущие 2);
$A \in \{0,\, 0{,}5,\, 1{,}0\}$ --- показатель регулярности по интервалу
между двумя последними занятиями ($\leq 3$ дней~--- 1{,}0;
$4$--$7$~--- 0{,}5; $> 7$~--- 0).

\textbf{Требуемый результат}~--- вектор рекомендуемых параметров:
\[
\mathbf{P}^{*}_e = (\mathbf{P}_e^{*\,\mathrm{load}},\, \mathbf{P}_e^{*\,\mathrm{sens}}),
\quad \mathbf{P}_e^{*\,\mathrm{sens}} = \mathbf{B}_e^{\mathrm{sens}}.
\]

Нагрузочная часть получается как корректировка назначения врача:
\[
\mathbf{P}_e^{*\,\mathrm{load}} = \mathrm{Adj}\bigl(\mathbf{B}_e^{\mathrm{load}},\, \delta\bigr),
\quad \delta \in \{-1,\, 0,\, +1\},
\]
где $\delta = -1$~--- упростить, $\delta = 0$~--- оставить назначение
врача, $\delta = +1$~--- осторожно усложнить.

Задача сводится к нахождению:
\[
\delta^{*} = R\bigl(\bar{S},\, T,\, A;\, \theta\bigr),
\qquad
\mathbf{P}_e^{*\,\mathrm{load}} = \mathrm{Adj}\bigl(\mathbf{B}_e^{\mathrm{load}},\, \delta^{*}\bigr).
\]

\subsubsection{Алгоритм решения}

Алгоритм состоит из пяти шагов.

\textbf{Шаг 1.} Загрузить $\mathbf{B}_e$ и $\mathcal{H}_{u,e}$
($N = 6$ последних сессий; если записей меньше~--- использовать все
доступные).

\textbf{Шаг 2.} Если число сессий $N < N_{\min} = 2$, вернуть
$\mathbf{P}^{*}_e = \mathbf{B}_e$ (рекомендация совпадает с назначением
врача).

\textbf{Шаг 3.} Вычислить $\bar{S}$, $T$, $A$ по формулам выше.
Дополнительно вычислить признак резкого спада:
$\mathrm{Drop} = 1$, если на двух последних сессиях подряд успешность
ниже, чем на третьей с конца, более чем на $15$~п.п.; иначе
$\mathrm{Drop} = 0$.

\textbf{Шаг 4.} Определить действие 

\[
\delta^{*} =
\left\{
\begin{array}{ll}
-1, & \text{если } \bar{S} < 60 \;\lor\; T < -10 \;\lor\; \mathrm{Drop} = 1 \;\lor\; A = 0, \\[6pt]
+1, & \text{если } \bar{S} \geq 80 \;\land\; T \geq 0 \;\land\; A \geq 0{,}5 \;\land\; \delta^{*} \neq -1, \\[6pt]
0, & \text{иначе.}
\end{array}
\right.
\]

Пороги выбраны по аналогии с пороговыми правилами из~\cite{noble2021,rule10}:
низкая успешность или ухудшение тренда~--- сигнал к упрощению;
устойчиво высокая успешность при регулярных занятиях~--- сигнал к
осторожному усложнению. За одну сессию допускается не более одного
шага сложности ($\delta \in \{-1, 0, +1\}$), что согласуется с
клиническим принципом постепенной прогрессии нагрузки~\cite{hall2016}.

\textbf{Шаг 5.} Вычислить
$\mathbf{P}_e^{*\,\mathrm{load}} = \mathrm{Adj}(\mathbf{B}_e^{\mathrm{load}},\, \delta^{*})$,
$\mathbf{P}_e^{*\,\mathrm{sens}} = \mathbf{B}_e^{\mathrm{sens}}$.

\subsubsection{Оператор $\mathrm{Adj}$}

Оператор $\mathrm{Adj}(\mathbf{B},\, \delta)$ применяет одно изменение
к одному нагрузочному параметру относительно $\mathbf{B}_e^{\mathrm{load}}$,
а не к параметрам последней сессии. При $\delta = 0$:
$\mathrm{Adj}(\mathbf{B},\, 0) = \mathbf{B}$. Все значения ограничиваются
допустимыми диапазонами интерфейса тренажёра.

\paragraph{Упражнения 1 и 2.}
Параметры: $n$ ($1 \ldots 20$)~--- количество объектов,
$\tau$ ($1 \ldots 30$~с)~--- время на объект.

При $\delta = -1$: если $\tau + 2 \leq 30$, то $\tau \leftarrow \tau + 2$;
иначе, если $n - 1 \geq 1$, то $n \leftarrow n - 1$.

При $\delta = +1$: если $\tau - 2 \geq 1$, то $\tau \leftarrow \tau - 2$;
иначе, если $n + 1 \leq 20$, то $n \leftarrow n + 1$.

\paragraph{Упражнение 3.}
Параметры: $n$ ($1 \ldots 20$), $v \in \{\text{Медленно},\,
\text{Средне},\,\text{Быстро}\}$.

При $\delta = -1$: понизить $v$ на один уровень; если $v =$ «Медленно»~---
уменьшить $n$ на~1.

При $\delta = +1$: повысить $v$ на один уровень; если $v =$ «Быстро»~---
увеличить $n$ на~1.

\paragraph{Упражнения 4, 5, 6.}
Параметры: $n$ ($1 \ldots 30$), $\tau$ ($10 \ldots 240$~с);
для 4 и~5 также $v$ (три уровня скорости).

При $\delta = -1$: для упр.~4, 5~--- понизить $v$; если скорость
минимальна (или упр.~6)~--- $\tau \leftarrow \min(\tau + 10,\, 240)$;
если $\tau$ максимально~--- $n \leftarrow \max(n - 1,\, 1)$.

При $\delta = +1$: для упр.~4, 5~--- повысить $v$; если скорость
максимальна (или упр.~6)~--- $\tau \leftarrow \max(\tau - 10,\, 10)$;
если $\tau$ минимально~--- $n \leftarrow \min(n + 1,\, 30)$.

\paragraph{Упражнение 7.}
Параметры: $n$, $\tau$, $r \in \{\text{Маленький},\,\text{Средний},\,
\text{Большой}\}$.

При $\delta = -1$: уменьшить $r$ на один уровень; иначе
$\tau \leftarrow \min(\tau + 10,\, 240)$; иначе $n \leftarrow n - 1$.

При $\delta = +1$: увеличить $r$ на один уровень; иначе
$\tau \leftarrow \max(\tau - 10,\, 10)$; иначе $n \leftarrow n + 1$.

\paragraph{Упражнения 8 и 9.}
Параметры: $n$, $\Delta t_c$ ($1 \ldots 10$~с)~--- интервал смены цвета,
$v$~--- скорость.

При $\delta = -1$: понизить $v$; иначе
$\Delta t_c \leftarrow \min(\Delta t_c + 1,\, 10)$; иначе $n \leftarrow n - 1$.

При $\delta = +1$: повысить $v$; иначе
$\Delta t_c \leftarrow \max(\Delta t_c - 1,\, 1)$; иначе $n \leftarrow n + 1$.

\subsubsection{Связь с анализом подходов}

Выбранный метод является частным случаем пороговой схемы
$a_k = R(x_k;\, \theta)$ из~\cite{noble2021,rule10,emg2020}, адаптированной
к специфике тренажёра:
\begin{itemize}
    \item опорная точка~--- назначение врача $\mathbf{B}_e$, а не
    параметры предыдущей сессии (как в~\cite{gomez2021});
    \item метрика $x_k$ заменена на вектор $(\bar{S}, T, A)$, вычисляемый
    по нескольким последним сессиям, а не по одному измерению;
    \item действие $\delta \in \{-1, 0, +1\}$ задаёт не только «менять /
    не менять», но и направление изменения;
    \item оператор $\mathrm{Adj}$ реализует конкретное изменение
    дискретных параметров с учётом ограничений интерфейса.
\end{itemize}

Нечёткая модель~\cite{gomez2021,castro2023} рассматривается как
возможное обобщение на следующем этапе: пороги $\theta$ могут быть
заменены функциями принадлежности без изменения общей структуры
алгоритма. Гибридный подход~\cite{zhou2024} и model-based методы
~\cite{noble2021,mi2020} избыточны при текущем объёме данных на
одного пациента.

\subsection{Реализация модуля рекомендаций в программном комплексе}

Модуль рекомендаций встроен в клиентское приложение тренажёра
(PyQt5, Python~3.10) и работает с базой данных MySQL через класс
\texttt{DatabasePreconnected} (модуль \texttt{database\_preconnected.py}).
Логика рекомендаций отделена от интерфейса и слоя доступа к данным:
вычисление признаков, принятие решения и корректировка параметров
сосредоточены в пакете \texttt{analytics/}, а главное окно
(\texttt{main.py}) отвечает только за вызов API и отображение
результата в элементах настройки упражнения.

\subsubsection{Архитектура пакета \texttt{analytics}}

Пакет состоит из четырёх модулей.

\begin{enumerate}
    \item \texttt{metrics.py}~--- расчёт вектора признаков
    $\boldsymbol{\Phi} = (\bar{S}, T, A)$ и признака резкого спада
    $\mathrm{Drop}$ по списку последних сессий;
    \item \texttt{policy.py}~--- функция $R$: по признакам возвращает
    $\delta \in \{-1, 0, +1\}$;
    \item \texttt{adj.py}~--- оператор $\mathrm{Adj}$: для каждого
    типа упражнения ($1 \ldots 9$) реализованы правила изменения
    нагрузочных параметров с учётом допустимых диапазонов интерфейса;
    \item \texttt{recommend.py}~--- точка входа
    \texttt{recommend(exercise\_id, baseline, history)}, объединяющая
    три предыдущих шага и возвращающая структуру результата
    (рекомендуемые параметры, $\delta$, признаки, флаг достаточности
    данных).
\end{enumerate}

Дополнительно модуль \texttt{ui\_binding.py} обеспечивает двустороннее
отображение параметров между словарём Python и виджетами формы
(\texttt{QSpinBox}, \texttt{QComboBox}) для каждого упражнения.
Такое разделение позволяет изменять правила в \texttt{policy.py}
и \texttt{adj.py} независимо от графического интерфейса.

\subsubsection{Хранение назначения врача}

Базовое назначение врача $\mathbf{B}_e$ хранится в таблице
\texttt{doctor\_baseline} (создаётся автоматически при первом
обращении). Структура таблицы:
\texttt{user\_id}, \texttt{exercise\_id}, \texttt{params\_json},
\texttt{updated\_at}; уникальный ключ по паре
(\texttt{user\_id}, \texttt{exercise\_id}).

Параметры сериализуются в JSON и содержат все настройки экрана
упражнения: нагрузочные (число объектов, время, скорость,
интервал смены цвета, диапазон шеи) и сенсорные (фон; для
упр.~1~--- также звук). Методы \texttt{save\_doctor\_baseline} и
\texttt{get\_doctor\_baseline} класса \texttt{DatabasePreconnected}
выполняют запись и чтение назначения.

Для расчёта рекомендаций используется метод
\texttt{get\_exercise\_sessions(user\_id, exercise\_id, limit=6)}:
он выбирает из таблицы результатов соответствующего упражнения
поля \texttt{apples\_count}, \texttt{caught\_apples},
\texttt{exercise\_date}, \texttt{coefficient}, \texttt{total\_score}
и возвращает записи в хронологическом порядке (от старых к новым).

\subsubsection{Интеграция в пользовательский интерфейс}

На экране настройки параметров упражнения (перед нажатием
«Начать упражнение») программно добавлены три кнопки:
\begin{itemize}
    \item «Сохранить как назначение врача»~--- считывает текущие
    значения виджетов формы и записывает их в \texttt{doctor\_baseline};
    \item «Рекомендуемые параметры»~--- загружает $\mathbf{B}_e$,
    последние 6~сессий, вызывает \texttt{recommend(...)} и подставляет
    полученные нагрузочные параметры в элементы управления;
    \item «Вернуть назначение врача»~--- восстанавливает сохранённое
    $\mathbf{B}_e$ без вызова модуля рекомендаций.
\end{itemize}

Сценарий работы:
\begin{enumerate}
    \item врач (или оператор) сохраняет назначение для упражнения;
    \item перед очередной тренировкой пациент нажимает
    «Рекомендуемые параметры»;
    \item система подставляет предложенные значения; пациент может
    начать упражнение, вернуть назначение врача или изменить
    параметры вручную.
\end{enumerate}

При недостатке истории ($N < 2$) модуль возвращает назначение врача
без изменений нагрузочной части. Сенсорные параметры (фон, звук)
при рекомендации не изменяются: на выход передаётся
$\mathbf{B}_e^{\mathrm{sens}}$ без модификации.

Таким образом, описанная математическая модель ($R$,
$\mathrm{Adj}$, привязка к $\mathbf{B}_e$) реализована в виде
отдельного модуля, интегрированного в программный комплекс
тренажёра и доступного пациенту на этапе настройки упражнения
перед тренировкой.
